\section{The Declaration Graph}
\label{sec:theorem-premise}

We now descend from modules to individual declarations. The import graph $G_{\mathrm{import}}$ of \S\ref{sec:import-graph} records which modules depend on which, treating each source file as an atomic unit; it cannot distinguish a module that contains one foundational definition from one that contains fifty routine lemmas. The \emph{declaration dependency graph} defined in this section resolves this limitation by recording the dependency of each individual declaration on its premises. At this finer granularity, the nodes are not files but theorems, definitions, and type constructors---the elementary units of mathematical content. Where the import graph captures the \emph{architecture} of the library, the declaration dependency graph captures the \emph{reasoning structure} of mathematics itself. In the next section (\S\ref{sec:namespace-graph}), we derive an intermediate family of graphs by aggregating declarations into namespaces; the declaration-level graph analyzed here is the ground truth from which those aggregations are computed.

The main findings of this section are threefold. First, the hub structure of the theorem--premise graph separates into two strata---a \emph{language infrastructure} layer (equality reasoning, numeral coercions) determined by the design of Lean's type system, and a \emph{mathematical infrastructure} layer (categories, real numbers, modules) determined by the logical structure of mathematics itself (\S\ref{sec:thm-types}--\S\ref{sec:thm-centrality}). Second, Louvain community detection recovers recognizable mathematical disciplines with moderate but imperfect alignment to the namespace hierarchy (NMI $= 0.34$), and only category theory achieves near-complete community--namespace correspondence (\S\ref{sec:thm-community}). Third, the cross-namespace density at the theorem level (50.9\%) far exceeds the module-level figure (37.1\%), revealing that mathematical reasoning violates cognitive classification more deeply than file organization alone suggests (\S\ref{sec:thm-cross-namespace}).

\subsection{Data and Definitions}
\label{sec:thm-data}

The theorem--premise graph is constructed from two complementary data sources. The \texttt{lean4export} tool~\cite{lean4export} extracts the full environment of a Lean~4 project---every declaration, its type, its proof term, and the constants appearing in each---into a machine-readable format. The \texttt{lean-training-data} tool~\cite{lean_training_data} extracts premise-level dependencies by analyzing which constants appear in the proof term of each declaration. We combine these extractions into a single dataset, published on HuggingFace as \texttt{Xinze-Li-Moqian/mathfactor}.

\paragraph{What is a premise?}
When a mathematician writes a proof in Lean, the tactic script is compiled into a \emph{proof term}---a fully explicit expression in Lean's dependent type theory. Every named constant from the environment that appears in this term is called a \textbf{premise} of the declaration. Consider a simple example:

\begin{codebox}
\textit{\color{gray}-- Tactic proof (what the mathematician writes):}\\
\textbf{theorem} \textcolor{red!70!black}{Nat.succ\_pos} (n : Nat) : 0 < n + 1 :=\\
\hspace{1.5em}\textcolor{red!70!black}{Nat.zero\_lt\_succ} n\\[3pt]
\textit{\color{gray}-- The proof term contains:}\\
\textit{\color{gray}--\hspace{0.5em}Nat.zero\_lt\_succ\hspace{1em}(named constant from environment $\to$ premise)}\\
\textit{\color{gray}--\hspace{0.5em}n\hspace{6.3em}(local variable, bound by theorem $\to$ NOT a premise)}
\end{codebox}

\noindent
The distinction is precise: a premise is an \emph{external} declaration---a theorem, definition, constructor, or other named constant already registered in the environment---not a local variable or bound parameter. In the example above, \decl{Nat.zero\_lt\_succ} is a premise because it is a theorem proved elsewhere; \texttt{n} is not, because it is a locally bound variable. In practice, the kernel records additional infrastructure premises (typeclass instances, notation abbreviations) that are invisible in the tactic script but present in the compiled term. It is these mechanically extracted premises that form the edges of $G_{\mathrm{thm}}$.

\begin{definition}[Declaration graph]
\label{def:thm-graph}
The \emph{declaration graph} is a directed graph $G_{\mathrm{thm}} = (\mathcal{D}, E_{\mathrm{thm}})$, where $\mathcal{D}$ is the set of all declarations---theorems, definitions, abbreviations, inductive types, constructors, opaques, quotients, and axioms---in \module{Mathlib}, and
\[
  E_{\mathrm{thm}} = \{(d_1, d_2) \in \mathcal{D} \times \mathcal{D} \mid d_2 \text{ appears as a premise in the proof term of } d_1\}.
\]
Each node carries two attributes: its \emph{kind} (one of eight declaration types) and its \emph{module} (the source file in which it is defined).
\end{definition}

The eight declaration kinds listed above are Lean~4's fundamental building blocks. Since they occupy distinct structural roles in $G_{\mathrm{thm}}$ (analyzed quantitatively in \S\ref{sec:thm-types}), we briefly describe each.

\paragraph{Theorem.}
A proved mathematical statement. The kernel stores the proof term for verification but erases it at runtime. Lean also accepts the keyword \texttt{lemma}, which is stored identically as a \texttt{theorem} in the compiled environment; our dataset does not distinguish the two.

\begin{codebox}
\textbf{theorem} \textcolor{red!70!black}{Nat.add\_comm} (a b : Nat) : a + b = b + a := \ldots
\end{codebox}

\noindent
In~$G_{\mathrm{thm}}$, a theorem's \emph{out-edges} point to the premises invoked in its proof, while its \emph{in-edges} come from later theorems that cite it. Theorems constitute $79.1\%$ of all declarations.

\paragraph{Definition.}
A computable function or constant whose body is retained at runtime. Definitions provide the computational content that theorems reason about.

\begin{codebox}
\textbf{def} \textcolor{red!70!black}{List.length} : List a -> Nat\\
\hspace{1.5em}| [] => 0 | \_ :: l => length l + 1
\end{codebox}

\noindent
Because many theorems invoke the same core operations (e.g., \decl{Set}, \decl{Eq.mpr}), definitions tend to accumulate high in-degree.

\paragraph{Abbreviation.}
A definitionally transparent shorthand, declared with \texttt{abbrev}. The elaborator always unfolds abbreviations, so they act as notational glue---typically wrapping typeclass projections or operator notation.

\begin{codebox}
\textbf{abbrev} \textcolor{red!70!black}{OfNat.ofNat} (n : Nat) [inst : OfNat a n] : a := inst.ofNat
\end{codebox}

\noindent
The most-cited node in the entire graph, \decl{OfNat.ofNat} (in-degree $89{,}936$), is an abbreviation: every declaration that mentions a numeric literal depends on it.

\paragraph{Inductive.}
The \texttt{inductive} keyword declares a new type by listing its introduction rules. The declaration itself becomes one node in $G_{\mathrm{thm}}$; each introduction rule becomes a separate \textbf{constructor} node (see below).

\begin{codebox}
\textbf{inductive} \textcolor{red!70!black}{Nat} \textbf{where}\\
\hspace{1.5em}| \textcolor{red!70!black}{zero} : Nat\\
\hspace{1.5em}| \textcolor{red!70!black}{succ} (n : Nat) : Nat
\end{codebox}

\noindent
Prominent inductive types in $G_{\mathrm{thm}}$ include \decl{CategoryTheory.Category} (in-degree $44{,}487$), \decl{Real} ($26{,}637$), and \decl{Module} ($23{,}404$).

\paragraph{Constructor.}
The introduction rules of an inductive type, auto-generated by the \texttt{inductive} declaration above. They are the \emph{only} way to build a value of the inductive type. For example, the declaration \texttt{inductive} \decl{Nat} above produces two constructor declarations:

\begin{codebox}
\textcolor{red!70!black}{Nat.zero} : Nat\hfill\textit{\color{gray}-- build zero}\\
\textcolor{red!70!black}{Nat.succ} : Nat -> Nat\hfill\textit{\color{gray}-- build a successor}
\end{codebox}

\noindent
Similarly, the inductive type \decl{Eq} (propositional equality) has a single constructor \decl{Eq.refl}, which is the sole way to introduce an equality proof. With in-degree $69{,}580$, it is the most-cited constructor in $G_{\mathrm{thm}}$.

\paragraph{Opaque.}
A definition whose body is sealed from the kernel: downstream proofs cannot unfold it. Opaques appear mainly in low-level system code such as floating-point arithmetic.

\begin{codebox}
\textbf{opaque} \textcolor{red!70!black}{Float.add} : Float -> Float -> Float
\end{codebox}

\paragraph{Quotient.}
Three kernel primitives that support quotient type constructions. Unlike all other declaration kinds, quotients are built into the Lean kernel rather than defined in library code:

\begin{codebox}
\textcolor{red!70!black}{Quot.mk}\hspace{1.1em}: (r : a -> a -> Prop) -> a -> Quot r\\
\textcolor{red!70!black}{Quot.lift}\hspace{0.3em}: (f : a -> b) -> \ldots\ -> Quot r -> b\\
\textcolor{red!70!black}{Quot.ind}\hspace{0.6em}: (p : Quot r -> Prop) -> \ldots\ -> (q : Quot r) -> p q
\end{codebox}

\paragraph{Axiom.}
A constant accepted without proof---the irreducible foundation. \module{Mathlib} uses exactly three axioms, on which all $308{,}129$ declarations ultimately rest:

\begin{codebox}
\textbf{axiom} \textcolor{red!70!black}{propext} : (a <-> b) -> a = b\hfill\textit{\color{gray}-- propositional extensionality}\\
\textbf{axiom} \textcolor{red!70!black}{Classical.choice} : Nonempty a -> a\hfill\textit{\color{gray}-- the axiom of choice}\\
\textbf{axiom} \textcolor{red!70!black}{Quot.sound} : r a b -> Quot.mk r a = Quot.mk r b
\end{codebox}

\paragraph{Edge directionality.}
Not all declaration kinds participate equally in $G_{\mathrm{thm}}$. An edge $(d_1, d_2)$ means that $d_1$'s proof or body \emph{cites} $d_2$; it does not imply an edge in the reverse direction. For example, the theorem \decl{Nat.lt\_irrefl} proves the irreflexivity of strict ordering on natural numbers: no $n$ satisfies $n < n$.

\begin{codebox}
\textit{\color{gray}-- Irreflexivity of < on natural numbers}\\
\textbf{theorem} \textcolor{red!70!black}{Nat.lt\_irrefl} (n : Nat) : Not (n < n) :=\\
\hspace{1.5em}\textcolor{red!70!black}{Nat.not\_succ\_le\_self} n
\end{codebox}

\noindent
The proof invokes \decl{Nat.not\_succ\_le\_self}, which establishes that $n + 1 \le n$ is impossible (since $<$ on \decl{Nat} is defined as $a < b \iff a + 1 \le b$, the statement $n < n$ reduces to $n + 1 \le n$). The kernel also records two infrastructure premises: \decl{LT.lt}, the abbreviation that provides the $<$~notation, and \decl{instLTNat}, the typeclass instance that connects $<$ to \decl{Nat}. This gives three edges:

\begin{figure}[H]
\centering
\begin{tikzpicture}[
  every node/.style={font=\small\ttfamily, inner sep=2pt, text=red!70!black},
  dot/.style={circle, fill=red!70!black, minimum size=3pt, inner sep=0pt},
  dep/.style={->, >=Stealth, thin, red!70!black},
]
  % Premises (cited) at top
  \node[dot, label={[align=center]above:{not\_succ\_le\_self\\[-2pt]{\normalfont\sffamily\scriptsize\color{gray}theorem}}}] (thm) at (-3.2,0) {};
  \node[dot, label={[align=center]above:{LT.lt\\[-2pt]{\normalfont\sffamily\scriptsize\color{gray}abbrev}}}]              (ab)  at (0,0) {};
  \node[dot, label={[align=center]above:{instLTNat\\[-2pt]{\normalfont\sffamily\scriptsize\color{gray}definition}}}]       (def) at (3.2,0) {};
  % Citing theorem at bottom
  \node[dot, label=below:{Nat.lt\_irrefl}] (t) at (0,-1.8) {};
  \draw[dep] (t) -- (thm);
  \draw[dep] (t) -- (ab);
  \draw[dep] (t) -- (def);
\end{tikzpicture}
\caption{Premise edges for the theorem \decl{Nat.lt\_irrefl}: the mathematical premise \decl{not\_succ\_le\_self} (a theorem) and two infrastructure premises \decl{LT.lt} (an abbreviation providing $<$ notation) and \decl{instLTNat} (a typeclass instance). The edges are directed: none of the three premises cites \decl{lt\_irrefl} in return.}
\label{fig:thm-single-edge}
\end{figure}

\noindent
In practice, the eight kinds arrange into a clear hierarchy of edge production and consumption:

\begin{table}[H]
\centering
\caption{Inter-kind edge flow in $G_{\mathrm{thm}}$. Each row is a source kind; each column is a target kind. Entries are thousands of edges; blanks denote fewer than $500$.}
\label{tab:inter-kind-flow}
\small
\begin{tabular}{l*{5}{r}}
\toprule
Source $\downarrow$ \textbackslash\ Target $\rightarrow$ & theorem & definition & abbrev & inductive & \textit{other} \\
\midrule
theorem      & $1{,}332$k & $2{,}498$k & $2{,}590$k & $835$k & $241$k \\
definition   & $18$k      & $285$k     & $284$k     & $163$k & $50$k \\
abbrev       & $2$k       & $17$k      & $26$k      & $20$k  & $3$k \\
constructor  & $1$k       & $12$k      & $20$k      & $16$k  & $6$k \\
inductive    &            & $2$k       & $3$k       & $6$k   & $5$k \\
axiom/opaque/quotient &   & $1$k       &            &        & \\
\bottomrule
\end{tabular}
\end{table}

\noindent
Theorems produce $89\%$ of all edges ($7.50$M of $8.44$M), primarily citing definitions and abbreviations---the computational and notational infrastructure on which proofs are built. Axioms and quotients are near-pure sinks: they receive citations but produce almost none (combined out-degree $= 3$). The graph is nearly acyclic; the only bidirectional edges ($4{,}713$ pairs) link inductive types to their own constructors (e.g., \decl{Nat}~$\leftrightarrow$~\decl{Nat.succ}), a structural artifact of how inductive definitions reference their introduction rules and vice versa.

\medskip
\noindent
The \decl{Nat.lt\_irrefl} example above showed a single theorem's premises; we now turn to a \emph{pair} of theorems to illustrate how shared premises create the hub structure that dominates $G_{\mathrm{thm}}$. Consider \decl{Nat.add\_left\_comm}, which proves $a + (b + c) = b + (a + c)$ for natural numbers. A tactic-level proof reads:

\begin{codebox}
\textbf{theorem} \textcolor{red!70!black}{Nat.add\_left\_comm} (a b c : Nat) :\\
\hspace{2em}a + (b + c) = b + (a + c) := \textbf{by}\\
\hspace{1.5em}rw [$\leftarrow$ \textcolor{red!70!black}{Nat.add\_assoc}, \textcolor{red!70!black}{Nat.add\_comm} a b, \textcolor{red!70!black}{Nat.add\_assoc}]
\end{codebox}

\noindent
The \texttt{lean-training-data} tool extracts premises not from the tactic script but from the compiled \emph{proof term}. For this theorem, the extracted premises are:

\smallskip
\begin{center}
\small
\begin{tabular}{ll}
\textit{Mathematical premises:} & \decl{Nat.add\_comm}, \decl{Nat.add\_assoc} \\
\textit{Equality infrastructure:} & \decl{Eq.refl}, \decl{Eq.symm}, \decl{Eq.mpr} \\
\textit{Typeclass instances:} & \decl{instHAdd}, \decl{HAdd.hAdd}, \decl{instAddNat}
\end{tabular}
\end{center}
\smallskip

\noindent
The first two are the mathematical lemmas that the proof logically depends on. The remaining six are \emph{language infrastructure}: equality combinators that the kernel requires to compose rewriting steps, and typeclass instances that elaborate the \texttt{+}~notation into the concrete function \decl{Nat.add}. Every edge in $E_{\mathrm{thm}}$ records one such dependency, so this single theorem contributes eight edges.

A closely related theorem, \decl{Nat.add\_right\_comm} ($a + b + c = a + c + b$), depends on \emph{exactly the same} eight premises. The resulting fragment of $G_{\mathrm{thm}}$ is shown in Figure~\ref{fig:thm-fragment}: both theorems share all three displayed premises, and \decl{Eq.refl}---a constructor with in-degree $69{,}580$---is the single most-cited declaration in the entire graph. This shared-hub pattern is the mechanism by which a small number of foundational declarations accumulate enormous in-degree.

\begin{figure}[H]
\centering
\begin{tikzpicture}[
  every node/.style={font=\small\ttfamily, inner sep=2pt, text=red!70!black},
  dot/.style={circle, fill=red!70!black, minimum size=3pt, inner sep=0pt},
  dep/.style={->, >=Stealth, thin, red!70!black},
]
  % Top row: shared premises (cited)
  \node[dot, label=above:{add\_comm}]  (ac) at (-3,0) {};
  \node[dot, label=above:{add\_assoc}] (aa) at (0,0)  {};
  \node[dot, label=above:{Eq.refl}]    (er) at (3,0)  {};

  % Bottom row: two citing theorems
  \node[dot, label=below:{add\_left\_comm}]  (alc) at (-2.2,-2) {};
  \node[dot, label=below:{add\_right\_comm}] (arc) at (2.2,-2)  {};

  % Edges from add_left_comm
  \draw[dep] (alc) -- (ac);
  \draw[dep] (alc) -- (aa);
  \draw[dep] (alc) -- (er);

  % Edges from add_right_comm
  \draw[dep] (arc) -- (ac);
  \draw[dep] (arc) -- (aa);
  \draw[dep] (arc) -- (er);
\end{tikzpicture}
\caption{A fragment of $G_{\mathrm{thm}}$: the theorems \decl{add\_left\_comm} and \decl{add\_right\_comm} both cite the premises \decl{add\_comm}, \decl{add\_assoc}, and \decl{Eq.refl}. Infrastructure premises (\decl{instHAdd}, \decl{Eq.symm}, etc.) are omitted for clarity. The shared node \decl{Eq.refl} (in-degree $69{,}580$) exemplifies the hub phenomenon analyzed in \S\ref{sec:thm-degree}. All names have the \texttt{Nat.}\ prefix removed except \decl{Eq.refl}.}
\label{fig:thm-fragment}
\end{figure}

\begin{remark}
The module \module{Algebra.Group.Defs} contains dozens of declarations. In $G_{\mathrm{import}}$, this entire module is a single vertex; in $G_{\mathrm{thm}}$, each declaration within it becomes a separate vertex, and each premise reference becomes a separate edge. A single import edge in $G_{\mathrm{import}}$ thus unfolds into potentially hundreds of premise edges in $G_{\mathrm{thm}}$---this is why $|E_{\mathrm{thm}}|$ exceeds $|E_{\mathrm{import}}|$ by a factor of $360$. This expansion is itself the transition from process (the human decision to bundle declarations into a module) to product (each logical dependency made explicit).
\end{remark}

At the commit analyzed, $|\mathcal{D}| = 308{,}129$ and $|E_{\mathrm{thm}}| = 8{,}436{,}366$. In comparison with the module-level import graph of \S\ref{sec:import-graph} ($|\mathcal{M}| = 7{,}563$ nodes, $|E_{\mathrm{import}}| = 23{,}570$ edges), the theorem--premise graph has roughly $40\times$ more nodes and $360\times$ more edges. The enormous increase in edge count reflects the fact that a single module import implicitly carries hundreds of individual premise dependencies; $G_{\mathrm{thm}}$ makes each of these dependencies explicit.

\begin{table}[ht]
\centering
\caption{Declaration type distribution in $G_{\mathrm{thm}}$.}
\label{tab:thm-kinds}
\begin{tabular}{lrr}
\toprule
Kind & Count & Percentage \\
\midrule
theorem     & $243{,}725$ & $79.1\%$ \\
definition  & $48{,}664$  & $15.8\%$ \\
abbrev      & $6{,}667$   & $2.2\%$ \\
constructor & $4{,}755$   & $1.5\%$ \\
inductive   & $3{,}813$   & $1.2\%$ \\
opaque      & $499$       & $0.2\%$ \\
quotient    & $3$         & ${<}0.1\%$ \\
axiom       & $3$         & ${<}0.1\%$ \\
\bottomrule
\end{tabular}
\end{table}

Theorems constitute nearly four-fifths of all declarations, with definitions accounting for the bulk of the remainder. The six axiom and quotient declarations at the bottom of the table---described in the paragraphs above---are few in number but structurally distinguished as the only pure sinks in $G_{\mathrm{thm}}$.

The graph is almost entirely connected: $308{,}054$ of $308{,}129$ nodes ($99.98\%$) belong to a single weakly connected component. The remaining $75$ nodes distribute across $71$ tiny components, of which $69$ are singletons. This near-perfect connectivity---stronger even than the module-level graph, which is exactly connected---reflects the deep logical interdependence of mathematical declarations.

\subsection{Declaration Types as Structural Strata}
\label{sec:thm-types}

The eight declaration kinds occupy distinct structural niches in $G_{\mathrm{thm}}$. Table~\ref{tab:thm-kind-degree} disaggregates the degree statistics by kind, revealing that the aggregate asymmetry of \S\ref{sec:thm-degree} arises from the superposition of qualitatively different populations.

\begin{table}[ht]
\centering
\caption{Degree statistics by declaration type. In-degree measures how often a declaration is cited; out-degree measures how many premises it invokes.}
\label{tab:thm-kind-degree}
\small
\begin{tabular}{lrrrrrrrrr}
\toprule
& & \multicolumn{4}{c}{In-degree} & \multicolumn{4}{c}{Out-degree} \\
\cmidrule(lr){3-6} \cmidrule(lr){7-10}
Kind & Count & Mean & Med.\ & Max & Std & Mean & Med.\ & Max & Std \\
\midrule
axiom       & $3$       & $7{,}823$ & $169$ & $23{,}226$ & $10{,}892$ & $0.7$ & $1$ & $1$ & $0.5$ \\
abbrev      & $6{,}667$ & $438$     & $11$  & $89{,}936$ & $3{,}033$  & $10.0$ & $7$ & $106$ & $10.8$ \\
inductive   & $3{,}813$ & $273$     & $31$  & $44{,}487$ & $1{,}541$  & $4.1$ & $3$ & $54$ & $4.4$ \\
constructor & $4{,}755$ & $59$      & $7$   & $69{,}580$ & $1{,}104$  & $11.7$ & $8$ & $127$ & $10.9$ \\
definition  & $48{,}664$& $58$      & $4$   & $62{,}942$ & $739$      & $16.4$ & $13$ & $146$ & $13.7$ \\
theorem     & $243{,}725$& $5.5$   & $1$   & $58{,}686$ & $237$      & $30.8$ & $21$ & $522$ & $31.2$ \\
quotient    & $3$       & $305$     & $226$ & $596$      & $213$      & $0.3$ & $0$ & $1$ & $0.5$ \\
opaque      & $499$     & $0.5$     & $0$   & $21$       & $1.4$      & $3.9$ & $3$ & $33$ & $3.1$ \\
\bottomrule
\end{tabular}
\end{table}

The pattern is systematic. Axioms sit at the absolute bottom of the dependency chain: they have the highest average in-degree ($7{,}823$) and near-zero out-degree. Among the three axioms, \decl{propext} (propositional extensionality) accounts for $23{,}226$ incoming edges---it is invoked in roughly one out of every thirteen declarations. \decl{Classical.choice} ($169$) and \decl{Quot.sound} ($73$) are invoked far less frequently, suggesting that the vast majority of \module{Mathlib}'s reasoning is constructive at the proof-term level, with classical logic concentrated in specific branches.

Inductive types have high in-degree (mean~$273$, median~$31$) and very low out-degree (mean~$4.1$). These are the structural skeletons of mathematics: type definitions that many theorems depend on but that themselves depend on little. Their in-degree top~$10$ reads as a catalogue of core mathematical concepts: \decl{CategoryTheory.Category} ($44{,}487$), \decl{Real} ($26{,}637$), \decl{TopologicalSpace} ($25{,}310$), \decl{CategoryTheory.Functor} ($24{,}295$), \decl{Module} ($23{,}404$), \decl{CommRing} ($20{,}837$).

Theorems occupy the opposite pole: low in-degree (mean~$5.5$, median~$1$) and the highest out-degree (mean~$30.8$, max~$522$). Theorems are \emph{consumers} of dependencies, not providers. The few high-in-degree theorems are not deep mathematical results but equality reasoning lemmas: \decl{Eq.trans} ($58{,}686$), \decl{of\_eq\_true} ($48{,}801$), \decl{congrFun'} ($45{,}516$), \decl{Eq.symm} ($42{,}542$). These are workhorses of Lean's elaborator and tactic framework---proof infrastructure, not mathematical content.

Abbreviations (\texttt{abbrev}) have the most extreme in-degree distribution among non-axiom types: mean~$438$, maximum~$89{,}936$ (\decl{OfNat.ofNat}). These are type-class coercion paths that the elaborator inserts automatically during type-checking. Their high citation counts reflect the machinery of the proof assistant, not the structure of mathematical reasoning.

\begin{remark}
\label{rem:two-layers}
The degree-by-kind analysis reveals that the hub structure of $G_{\mathrm{thm}}$ decomposes into two distinct infrastructure layers:
\begin{enumerate}
\item \textbf{Language infrastructure.} The highest-cited declarations---\decl{OfNat.ofNat} ($89{,}936$), \decl{Eq.refl} ($69{,}580$), \decl{DFunLike.coe} ($65{,}437$), \decl{Eq.mpr} ($62{,}942$)---are artifacts of Lean's type system: numeral coercions, equality constructors, function-like coercions. Their prominence is determined by the design of the proof assistant, not by the logical structure of mathematics. These nodes are invisible in $G_{\mathrm{import}}$ (which treats entire modules as atoms) and constitute the first layer of process-trace in the product.

\item \textbf{Mathematical infrastructure.} The highest-cited \emph{inductive types}---\decl{CategoryTheory.Category}, \decl{Real}, \decl{TopologicalSpace}, \decl{Module}, \decl{CommRing}---are core mathematical concepts. Their prominence reflects the logical structure of the formalized mathematics itself: category theory, analysis, and algebra are genuinely foundational in the dependency sense. This layer is the product's own structure, free of proof-assistant artifacts.
\end{enumerate}
The import graph of \S\ref{sec:import-graph} conflates these two layers (a module contains both language infrastructure and mathematical content). The theorem--premise graph separates them, enabling the analyst to distinguish what belongs to mathematics from what belongs to its medium of expression.
\end{remark}

\subsection{Cross-Namespace Flow}
\label{sec:thm-cross-namespace}

We decompose the $8{,}436{,}366$ edges of $G_{\mathrm{thm}}$ into three tiers based on the relationship between the source and target declarations' names. Here ``module'' is the source file containing the declaration, and ``namespace'' is the first component of its fully qualified name (e.g., the namespace of \decl{Nat.lt\_irrefl} is \ns{Nat})---a declaration-name prefix, distinct from the directory-based $\mathrm{dir}(m)$ of \S\ref{sec:cross-namespace}.

\begin{table}[ht]
\centering
\caption{Edge classification in $G_{\mathrm{thm}}$ by module and namespace relationship.}
\label{tab:thm-edge-breakdown}
\begin{tabular}{lrr}
\toprule
Category & Edges & Percentage \\
\midrule
Same module                    & $811{,}649$   & $9.6\%$ \\
Cross-module, same namespace   & $805{,}772$   & $9.6\%$ \\
Cross-namespace                & $4{,}296{,}412$ & $50.9\%$ \\
Missing module info            & $2{,}522{,}533$ & $29.9\%$ \\
\bottomrule
\end{tabular}
\end{table}

Only $9.6\%$ of edges remain within the same module file, and another $9.6\%$ cross module boundaries while staying within the same top-level namespace. The majority---$50.9\%$---cross namespace boundaries entirely. The $29.9\%$ ``missing'' category corresponds to declarations in the root namespace (lacking a module prefix), which predominantly involve core Lean infrastructure (\decl{Eq}, \decl{Iff}, \decl{OfNat}, etc.).\footnote{If root-level declarations are treated as their own pseudo-namespace, this category would further increase the cross-namespace fraction, since root infrastructure is cited by declarations across all namespaces.}

The comparison with the module-level cross-namespace rate is stark. In $G_{\mathrm{import}}^{-}$, $62.9\%$ of edges stay within the same namespace (\S\ref{sec:cross-namespace}); in $G_{\mathrm{thm}}$, at most $19.2\%$ do ($9.6\% + 9.6\%$). This threefold increase in cross-namespace density arises because module imports package many individual dependencies into a single edge: a module that imports another module within the same namespace may, at the theorem level, depend on hundreds of declarations from \emph{other} namespaces that were transitively imported. The module boundary acts as a cognitive container that conceals the true extent of cross-disciplinary dependency.

\begin{table}[ht]
\centering
\caption{Top~10 cross-namespace edge pairs in $G_{\mathrm{thm}}$, ordered by edge count.}
\label{tab:thm-cross-ns-pairs}
\small
\begin{tabular}{llr}
\toprule
Namespace A & Namespace B & Edges \\
\midrule
\ns{AlgebraicGeometry}  & \ns{CategoryTheory}    & $38{,}042$ \\
\ns{CategoryTheory}     & \ns{Quiver}            & $29{,}286$ \\
\ns{CategoryTheory}     & \ns{Eq}                & $27{,}927$ \\
\ns{CategoryTheory}     & \ns{HomologicalComplex} & $15{,}829$ \\
\ns{MeasureTheory}      & \ns{Real}              & $15{,}668$ \\
\ns{MeasureTheory}      & \ns{Set}               & $11{,}954$ \\
\ns{MeasureTheory}      & \ns{ProbabilityTheory} & $11{,}734$ \\
\ns{Eq}                 & \ns{MeasureTheory}     & $10{,}385$ \\
\ns{Complex}            & \ns{Real}              & $10{,}370$ \\
\ns{ENNReal}            & \ns{MeasureTheory}     & $10{,}302$ \\
\bottomrule
\end{tabular}
\end{table}

The strongest cross-namespace coupling is \ns{AlgebraicGeometry} $\leftrightarrow$ \ns{CategoryTheory} ($38{,}042$ edges), reflecting algebraic geometry's fundamental dependence on categorical language---schemes, sheaves, and functors are inherently categorical constructions. The \ns{MeasureTheory} $\leftrightarrow$ \ns{Real} coupling ($15{,}668$ edges) represents the bridge between abstract measure theory and concrete analysis on $\mathbb{R}$. The presence of \ns{Eq} as a namespace in multiple top pairs confirms the pervasive role of equality reasoning infrastructure across all mathematical domains.

\subsection{The Leaf Theorem Phenomenon}
\label{sec:thm-leaves}

Of the $243{,}725$ theorems in $G_{\mathrm{thm}}$, $109{,}088$ ($44.8\%$) have in-degree zero: they are never cited by any other declaration. These \emph{leaf theorems} form the outermost surface of the library---the boundary where formalized knowledge meets the external world.

\begin{table}[ht]
\centering
\caption{Zero-citation rates by namespace (extremes among namespaces with $\ge 100$ theorems).}
\label{tab:thm-leaves}
\small
\begin{tabular}{lrrrlrrr}
\toprule
\multicolumn{4}{c}{Highest zero-citation rate} & \multicolumn{4}{c}{Lowest zero-citation rate} \\
\cmidrule(lr){1-4} \cmidrule(lr){5-8}
Namespace & Total & Zero & Rate & Namespace & Total & Zero & Rate \\
\midrule
\ns{AddEquiv}           & $129$  & $119$ & $92.2\%$ & \ns{Primrec}            & $133$ & $17$ & $12.8\%$ \\
\ns{EMetric}            & $192$  & $163$ & $84.9\%$ & \ns{AnalyticAt}         & $102$ & $15$ & $14.7\%$ \\
\ns{DomMulAct}          & $105$  & $89$  & $84.8\%$ & \ns{AkraBazziRecurrence}& $129$ & $21$ & $16.3\%$ \\
\ns{Ordering}           & $127$  & $105$ & $82.7\%$ & \ns{ContinuousWithinAt} & $119$ & $20$ & $16.8\%$ \\
\ns{LocallyConstant}    & $113$  & $93$  & $82.3\%$ & \ns{HasFDerivWithinAt}  & $120$ & $21$ & $17.5\%$ \\
\ns{MulEquiv}           & $243$  & $194$ & $79.8\%$ & \ns{ConvexOn}           & $111$ & $20$ & $18.0\%$ \\
\ns{AddSubgroup}        & $111$  & $88$  & $79.3\%$ & \ns{LipschitzWith}      & $113$ & $21$ & $18.6\%$ \\
\ns{AddMonoidHom}       & $138$  & $108$ & $78.3\%$ & \ns{DifferentiableAt}   & $113$ & $21$ & $18.6\%$ \\
\bottomrule
\end{tabular}
\end{table}

The namespaces with the highest zero-citation rates---\ns{AddEquiv} ($92.2\%$), \ns{EMetric} ($84.9\%$), \ns{MulEquiv} ($79.8\%$), \ns{AddSubgroup} ($79.3\%$), \ns{AddMonoidHom} ($78.3\%$)---share a common pattern: they are generated by the \texttt{@[to\_additive]} attribute, which systematically creates additive counterparts of multiplicative lemmas. These ``mirrored'' theorems exist for API completeness---ensuring that every statement about groups has an analogous statement about additive groups---but are rarely invoked directly in proof terms, because the original multiplicative version or a more general statement is typically preferred.

The namespaces with the lowest zero-citation rates tell a complementary story. \ns{Primrec} ($12.8\%$), \ns{AnalyticAt} ($14.7\%$), and the differentiation predicates (\ns{HasFDerivWithinAt}, \ns{DifferentiableAt}, \ns{ContinuousWithinAt}) produce lemmas that are heavily reused as building blocks. These are the workbench tools of analysis and computability theory---nearly every theorem in these namespaces feeds into downstream proofs.

\begin{remark}
\label{rem:leaf-theorems}
The $44.8\%$ zero-citation rate means that nearly half of all theorems in $G_{\mathrm{thm}}$ are leaves. These leaves constitute the \emph{API surface} of \module{Mathlib}---the interface through which the library's knowledge is consumed by external users (mathematicians applying results, AI systems selecting premises, downstream projects building on Mathlib). Their existence is not waste but function: a library that formalized only internally-reused lemmas would be useless to its consumers.

The dichotomy between high-zero-rate namespaces (systematically generated mirrored API) and low-zero-rate namespaces (hand-crafted building blocks) reveals a second process-trace in the product. The \texttt{@[to\_additive]} machinery leaves a distinctive structural fingerprint---entire subtrees of leaf theorems that inflate the graph's surface area without deepening its logical content. This automation artifact is invisible at the module level (a module contains both original and mirrored theorems) but becomes visible at the theorem level, where individual declarations can be classified by their provenance.
\end{remark}

\subsection{Analysis}
\label{sec:thm-conclusion}

The theorem--premise graph $G_{\mathrm{thm}}$ is a more faithful record of mathematical reasoning than the import graph $G_{\mathrm{import}}$. It dissolves the module boundary, exposing the individual logical dependencies that file-level analysis aggregates away. Our analysis of $G_{\mathrm{thm}}$ yields three central findings.

\medskip
\noindent\textit{1. The double infrastructure.}\\
The hub structure of $G_{\mathrm{thm}}$ decomposes into two layers that the import graph conflates: a \emph{language infrastructure} layer (\decl{OfNat.ofNat}, \decl{Eq.refl}, \decl{DFunLike.coe}) determined by Lean's type system, and a \emph{mathematical infrastructure} layer (\decl{CategoryTheory.Category}, \decl{Real}, \decl{TopologicalSpace}) determined by the logical structure of formalized mathematics (Remark~\ref{rem:two-layers}). This separation is invisible at the module level, where both layers cohabit the same source files. The language layer is a trace of the process---the specific proof assistant chosen---embedded in the product.

\medskip
\noindent\textit{2. Disciplinary autonomy as emergent community structure.}\\
Louvain community detection recovers seven major mathematical disciplines from the dependency topology alone (Table~\ref{tab:thm-communities}), with a modularity of $0.48$. The alignment with namespace organization, however, is only moderate (NMI $= 0.34$, ARI $= 0.18$; Table~\ref{tab:thm-nmi}). Among the major namespaces, only \ns{CategoryTheory} achieves near-complete community correspondence ($97.4\%$ concentration in a single community), suggesting that it is a genuinely autonomous mathematical discipline within \module{Mathlib}---self-contained in its dependency structure to a degree that no other branch achieves. All other communities are cross-namespace mixtures, confirming that the dependency topology discovers a structure related to, but distinct from, the human-imposed namespace hierarchy.

\medskip
\noindent\textit{3. The depth of cross-disciplinary entanglement.}\\
The cross-namespace edge density of $G_{\mathrm{thm}}$ ($50.9\%$; Table~\ref{tab:thm-edge-breakdown}) far exceeds that of $G_{\mathrm{import}}^{-}$ ($37.1\%$; \S\ref{sec:cross-namespace}). This gap reveals that module-level analysis \emph{underestimates} the true extent of mathematical interdependence. Module boundaries bundle cross-namespace theorem dependencies into single intra-directory import edges, creating an illusion of disciplinary containment that dissolves at finer granularity.

\medskip

Table~\ref{tab:thm-vs-module} summarizes the key quantitative contrasts between the two levels of analysis.

\begin{table}[ht]
\centering
\caption{Module-level ($G_{\mathrm{import}}$) vs.\ theorem-level ($G_{\mathrm{thm}}$) comparison.}
\label{tab:thm-vs-module}
\begin{tabular}{lrr}
\toprule
Metric & $G_{\mathrm{import}}$ / $G_{\mathrm{import}}^{-}$ & $G_{\mathrm{thm}}$ \\
\midrule
Nodes                        & $7{,}563$          & $308{,}129$ \\
Edges                        & $23{,}570$ / $19{,}448$ & $8{,}436{,}366$ \\
Cross-namespace edges         & $37.1\%$           & $50.9\%$ \\
In-degree max                & $167$              & $89{,}936$ \\
Community--namespace NMI     & ---                & $0.34$ \\
Single-node removal impact   & $29$ components    & $\le 6$ nodes disconnected \\
Zero-citation nodes          & $1{,}433$ sources  & $109{,}088$ leaf theorems ($44.8\%$) \\
\bottomrule
\end{tabular}
\end{table}

We close by connecting these findings to the central tension of the paper. The theorem--premise graph is a more direct projection of the logical product than the import graph: its hubs, communities, and cross-namespace flows are less contaminated by file-organization conventions. Yet even at this finer resolution, the process leaves indelible traces in the product. The language infrastructure layer (Remark~\ref{rem:two-layers}) records the design choices of Lean's type system---a different proof assistant would produce a different in-degree ranking, even for the same mathematics. The $44.8\%$ leaf theorem rate (Remark~\ref{rem:leaf-theorems}) includes a substantial fraction of \texttt{@[to\_additive]} mirrors---a process-level automation that inflates the graph's surface without deepening its logical content. The out-degree distribution's lognormal shape (Remark~\ref{rem:thm-degree-asymmetry}) encodes a cognitive constraint: the bounded complexity of proofs that humans can construct or that tactics can synthesize. Even in the most product-like graph we have analyzed, the process is not erased; it is encoded.

For AI premise selection systems, the two-layer hub structure suggests a practical preprocessing step: filtering language-infrastructure premises (which carry no mathematical content) from training data may improve the signal-to-noise ratio. The theorem/lemma distinction, recoverable from source code (\S\ref{sec:thm-vs-lemma}), provides an additional human-annotated importance signal that could inform premise ranking.
